% ============================================================================
%  Closure Test of the Cell-Energy Reweighting Method
%  for Photon Shower Shape Corrections in MC23e
% ============================================================================
\documentclass[12pt,a4paper]{article}

% --- Packages ---------------------------------------------------------------
\usepackage[utf8]{inputenc}
\usepackage[T1]{fontenc}
\usepackage{lmodern}
\usepackage[margin=2.5cm,headheight=15pt]{geometry}
\usepackage{amsmath,amssymb}
\usepackage{graphicx}
\usepackage{pdfpages}
\usepackage{booktabs}
\usepackage{caption}
\usepackage{subcaption}
\usepackage{hyperref}
\usepackage{cleveref}
\usepackage{xcolor}
\usepackage{float}
\usepackage{siunitx}
\usepackage{setspace}
\usepackage{fancyhdr}

% --- Page style -------------------------------------------------------------
\pagestyle{fancy}
\fancyhf{}
\fancyhead[L]{\small Closure Test — Cell-Energy Reweighting}
\fancyhead[R]{\small ATLAS Internal}
\fancyfoot[C]{\thepage}
\renewcommand{\headrulewidth}{0.4pt}

% --- Graphics path ----------------------------------------------------------
\graphicspath{
  {../output/closure_test_5pct/}
  {../output/closure_test_3pct/}
  {../output/closure_test_1pct/}
}

% --- Hyperref setup ---------------------------------------------------------
\hypersetup{
  colorlinks=true,
  linkcolor=blue!60!black,
  citecolor=blue!60!black,
  urlcolor=blue!60!black
}

% --- Title ------------------------------------------------------------------
\title{%
  \textbf{Closure Test of the Cell-Energy Reweighting Method\\[4pt]
  for Photon Shower Shape Corrections in MC23e}
}
\author{ATLAS Internal}
\date{February 2026}

% ============================================================================
\begin{document}
\maketitle
\thispagestyle{fancy}

\begin{abstract}
\noindent
A closure test of the cell-energy reweighting method for correcting photon
shower shape variables in the ATLAS electromagnetic calorimeter is presented.
The method operates at the level of individual calorimeter cell energies in the
second sampling layer, correcting all derived shower shape variables
simultaneously. Pseudo-data are generated from MC23e $Z\to ee\gamma$
simulation by applying controlled per-cell energy distortions, and two
correction methods — flat cell reweighting
(§5.1 of ATL-COM-PHYS-2021-640) and energy-dependent cell reweighting
(§5.2) — are derived and validated. Using the full sample of 9.6~million
events, the energy-dependent cell reweighting method achieves $\chi^2/n_\text{df}$ reductions
of 60--95\% relative to uncorrected Monte Carlo across all pseudorapidity bins
and shower shape variables, demonstrating excellent closure.
\end{abstract}

\tableofcontents
\newpage

% ============================================================================
\section{Introduction}
\label{sec:introduction}
% ============================================================================

Photon identification in the ATLAS experiment relies heavily on
electromagnetic calorimeter shower shape variables, which characterise the
lateral energy distribution of the shower in the second sampling layer of the
liquid argon calorimeter. Discrepancies between data and Monte Carlo (MC)
simulation in these variables can lead to systematic biases in photon
identification efficiency measurements and must be corrected.

The cell-energy reweighting method, described in
ATL-COM-PHYS-2021-640~(Section~5), provides a data-driven correction by
adjusting the energy of individual calorimeter cells in the
$7\times 11$~($\eta\times\phi$) cluster of the second electromagnetic layer.
Rather than correcting shower shape variables directly, this approach modifies
the underlying cell-level energy fractions, thereby simultaneously correcting
all derived quantities in a physically consistent manner.

This report presents a closure test of the cell-energy reweighting method
applied to MC23e $Z\to ee\gamma$ Monte Carlo samples. The purpose of the
closure test is to validate the correction procedure in a controlled
environment where the ``true'' answer is known by construction. Pseudo-data
are generated from MC by applying known per-cell distortions that mimic
typical data--MC differences, corrections are derived and applied using two
complementary methods, and the corrected distributions are compared to the
pseudo-data target to quantify the degree of closure.

Two shower shape variables are examined:
\begin{itemize}
  \item $R_\eta = E(3\times7)\,/\,E(7\times7)$, the ratio of energies in
        the $3\eta\times 7\phi$ and $7\eta\times 7\phi$ sub-windows;
  \item $R_\phi = E(3\times3)\,/\,E(3\times7)$, the ratio of energies in
        the $3\eta\times 3\phi$ and $3\eta\times 7\phi$ sub-windows.
\end{itemize}

% ============================================================================
\section{Data and Inputs}
\label{sec:data}
% ============================================================================

% ---------- 2.1 MC sample --------------------------------------------------
\subsection{Monte Carlo sample}
\label{sec:data:mc}

The input sample is a MC23e $Z\to ee\gamma$ simulation (DSID~700770)
produced with the ATLAS Run~3 simulation campaign. The sample contains
approximately 9.6~million events.

Each event provides the cell-level energy deposits for the photon candidate
in the second electromagnetic calorimeter layer, stored as a vector of 77
values corresponding to the cells of the $7\times 11$~($\eta\times\phi$)
cluster. The cell energies are given in MeV. The pseudorapidity of the
cluster barycentre in the second layer, $\eta_2$, is also available and is
used to assign each photon to one of 14~$|\eta|$ bins.

% ---------- 2.2 Eta binning ------------------------------------------------
\subsection{Pseudorapidity binning}
\label{sec:data:eta}

The analysis is performed differentially in $|\eta|$, using 14~bins that
follow the standard ATLAS electromagnetic calorimeter segmentation:
\begin{equation}
  |\eta| \in
  \bigl\{
    [0,0.2),\;
    [0.2,0.4),\;
    [0.4,0.6),\;
    [0.6,0.8),\;
    [0.8,1.0),\;
    [1.0,1.2),\;
    [1.2,1.3),\;
    [1.3,1.37),
  \bigr.
  \label{eq:etabins1}
\end{equation}
\begin{equation}
  \bigl.
    [1.52,1.6),\;
    [1.6,1.8),\;
    [1.8,2.0),\;
    [2.0,2.2),\;
    [2.2,2.4)
  \bigr\}.
  \label{eq:etabins2}
\end{equation}
The gap between bins~7 ($|\eta|<1.37$) and~9 ($|\eta|\geq 1.52$)
corresponds to the barrel--endcap transition region of the electromagnetic
calorimeter, where the detector geometry changes and photon reconstruction is
not performed. This region is excluded from the analysis.

% ---------- 2.3 Quality selection -------------------------------------------
\subsection{Cluster quality selection}
\label{sec:data:quality}

The MC23e sample contains all photon candidates from the standard
reconstruction, without tight identification or high-$p_\text{T}$ selection.
A non-negligible fraction of low-quality clusters is therefore present. Two
quality requirements are applied:

\paragraph{Non-negative cell energies.}
Cell energies can take negative values due to electronic noise subtraction in
the calorimeter readout. Negative energies produce unphysical cell fractions
that distort the mean fraction calculations and the profile corrections. All
negative cell energies are set to zero prior to further processing.

\paragraph{Central cell dominance.}
Well-reconstructed electromagnetic showers have the highest energy deposition
in the central cell of the cluster (cell at row~3, column~5 in the
$7\times 11$ grid). Clusters where a peripheral cell has a larger energy
fraction than the central cell are indicative of mis-centred or pathological
clusters and are rejected.

These quality requirements are applied consistently across all stages of the
analysis.

% ---------- 2.4 Fraction binning --------------------------------------------
\subsection{Energy fraction binning}
\label{sec:data:fracbins}

The energy-dependent cell reweighting (Method~2, described below) uses profile
histograms to record the mean correction as a function of the MC cell energy
fraction. The profile binning uses 12~bins spanning the range $[-1,1]$:
\begin{equation}
  f_k \in
  \{-1,\;-0.5,\;-0.4,\;-0.3,\;-0.2,\;-0.1,\;0,\;0.1,\;0.2,\;0.3,\;0.4,\;0.5,\;1\}.
  \label{eq:fracbins}
\end{equation}
This binning provides six bins of width~0.1 in the region $[0,0.5]$ where
the physical cell fractions are concentrated, ensuring adequate resolution for
the energy-dependent cell reweighting.

% ============================================================================
\section{Methodology}
\label{sec:method}
% ============================================================================

% ---------- 3.1 Overview ----------------------------------------------------
\subsection{Overview}
\label{sec:method:overview}

The closure test proceeds in four stages:
\begin{enumerate}
  \item \textbf{Pseudo-data generation}: MC cell energies are distorted by a
        known transformation to produce pseudo-data that serves as the
        correction target.
  \item \textbf{Correction derivation}: The pseudo-data and original MC are
        compared cell-by-cell to extract correction parameters.
  \item \textbf{Correction application}: The derived corrections are applied
        to the original MC to produce reweighted MC.
  \item \textbf{Validation}: Shower shape distributions computed from the
        reweighted MC are compared to the pseudo-data distributions using a
        $\chi^2$ metric.
\end{enumerate}
Two correction methods are implemented and compared, corresponding to
Sections~5.1 and~5.2 of ATL-COM-PHYS-2021-640.

% ---------- 3.2 Pseudo-data generation --------------------------------------
\subsection{Pseudo-data generation}
\label{sec:method:pseudodata}

The distortion model is designed to mimic the qualitative features of typical
data--MC discrepancies in photon shower shapes: MC simulation tends to produce
narrower showers than data, corresponding to an excess of energy in the
cluster centre relative to the periphery. Fudge factors correct this by
broadening MC showers. The pseudo-data therefore simulates data-like broader
showers by redistributing energy from the central cells toward the tails.

For each event passing the quality selection, the energy of cell~$k$ is
modified according to:
\begin{equation}
  e'_k \;=\; e_k \,\cdot\, \bigl(1 + \delta\cdot b_k + \varepsilon_k\bigr),
  \label{eq:distortion}
\end{equation}
where:
\begin{itemize}
  \item $\delta$ is the \textbf{systematic distortion level}, a global
        parameter controlling the magnitude of the bias. The closure test is
        performed at three values: $\delta=0.01$~(1\%), $\delta=0.03$~(3\%),
        and $\delta=0.05$~(5\%).
  \item $b_k$ is a \textbf{fixed bias pattern} determined by the Chebyshev
        distance~$d_k$ of cell~$k$ from the central cell~$(3,5)$:
        \begin{equation}
          b_k = \begin{cases}
            -1.0 & d_k = 0 \;\text{(central cell)}, \\
            -0.4 & d_k = 1 \;\text{(first ring)}, \\
            +0.4 & d_k = 2 \;\text{(second ring)}, \\
            +1.0 & d_k \geq 3 \;\text{(outer cells)}.
          \end{cases}
          \label{eq:biaspattern}
        \end{equation}
        This pattern suppresses the central core and enhances the tails,
        producing a broadening of the shower that mimics the standard
        ATLAS data--MC difference.
  \item $\varepsilon_k \sim \mathcal{N}(0,\,\sigma)$ is a
        \textbf{per-event, per-cell Gaussian random noise} with
        $\sigma = \delta/2$, representing stochastic event-by-event
        fluctuations that cannot be corrected on a per-event basis but
        contribute to the overall distribution shape.
\end{itemize}

After the per-cell scaling, total cluster energy is conserved by rescaling:
\begin{equation}
  e'_k \;\leftarrow\; e'_k \cdot \frac{E_\text{tot}}{\sum_j e'_j},
  \qquad E_\text{tot} = \sum_j e_j.
  \label{eq:energycons}
\end{equation}
This ensures that the distortion affects only the lateral energy distribution
and not the overall energy scale.

Events that do not pass the quality selection are passed through unchanged. A
fixed random seed is used to ensure full reproducibility.

% ---------- 3.3 Correction derivation ---------------------------------------
\subsection{Correction derivation}
\label{sec:method:derive}

For each event, cell energy fractions are computed as:
\begin{equation}
  f_k = \frac{e_k}{E_\text{tot}} = \frac{e_k}{\sum_j e_j},
  \label{eq:fraction}
\end{equation}
separately for the MC sample ($f_k^\text{MC}$) and the pseudo-data sample
($f_k^\text{PS}$).

\paragraph{Method~1: Flat cell reweighting (§5.1).}
This method derives a single additive correction per cell per $|\eta|$ bin:
\begin{equation}
  \Delta_k^{(n)} =
  \langle f_k^\text{PS}\rangle_n
  - \langle f_k^\text{MC}\rangle_n,
  \label{eq:m1derive}
\end{equation}
where $\langle\cdot\rangle_n$ denotes the mean over all events in $|\eta|$
bin~$n$. The correction captures the average systematic difference between MC
and pseudo-data energy fractions but is independent of the event-by-event
value of~$f_k$.

\paragraph{Method~2: Energy-dependent cell reweighting (§5.2).}
This method derives an energy-dependent correction that varies as a function
of the MC cell fraction. For each cell~$k$ and $|\eta|$ bin~$n$, a profile
histogram records:
\begin{equation}
  \alpha_k^{(n)}(f) =
  \bigl\langle
    f_k^\text{PS} - f_k^\text{MC}
  \bigr\rangle \Big|_{f_k^\text{MC} = f},
  \label{eq:m2derive}
\end{equation}
i.e.\ the mean difference between pseudo-data and MC fractions, computed in
bins of the MC fraction~$f_k^\text{MC}$. This captures the fact that the
correction may depend on the cell's energy share: cells carrying a larger
fraction of the total energy may require a different correction than cells
with a smaller fraction.

One profile histogram is created for each of the 77~cells in each of the
14~$|\eta|$ bins, giving $77\times 14 = 1078$ profiles in total.

% ---------- 3.4 Correction application --------------------------------------
\subsection{Correction application}
\label{sec:method:apply}

The corrections are applied to the original MC by modifying the cell energy
fractions and then reconstructing the cell energies.

\textbf{Method~1} (flat cell reweighting, §5.1) applies a flat shift:
\begin{equation}
  f_k^\text{corr} = f_k^\text{MC} + \Delta_k^{(n)}.
  \label{eq:m1apply}
\end{equation}

\textbf{Method~2} (energy-dependent cell reweighting, §5.2) applies a
correction obtained by interpolating the profile histogram at the value of
the MC fraction:
\begin{equation}
  f_k^\text{corr} = f_k^\text{MC}
    + \alpha_k^{(n)}\!\bigl(f_k^\text{MC}\bigr).
  \label{eq:m2apply}
\end{equation}

In both cases, the corrected cell energies are reconstructed with energy
conservation:
\begin{equation}
  e_k^\text{corr} = \frac{f_k^\text{corr}}{\sum_j f_j^\text{corr}}
  \cdot E_\text{tot}.
  \label{eq:correnergycons}
\end{equation}
Events failing the quality selection are passed through without correction.

% ---------- 3.5 Shower shape computation ------------------------------------
\subsection{Shower shape computation}
\label{sec:method:shapes}

After the correction, the two shower shape variables are recomputed from
the modified cell energies.

The energy in a rectangular sub-window of dimensions $m\times n$
($\eta\times\phi$) centred on the cluster core is defined as:
\begin{equation}
  E(m\times n) =
  \sum_{\substack{
    i_\eta = c_\eta - \lfloor m/2\rfloor \\
    \phantom{i_\eta = c_\eta}\;\;\ldots\;\;c_\eta + \lfloor m/2\rfloor
  }}
  \;\;
  \sum_{\substack{
    i_\phi = c_\phi - \lfloor n/2\rfloor \\
    \phantom{i_\phi = c_\phi}\;\;\ldots\;\;c_\phi + \lfloor n/2\rfloor
  }}
  e_{i_\phi + 11 \cdot i_\eta},
  \label{eq:subwindow}
\end{equation}
where $(c_\eta, c_\phi) = (3,5)$ is the central cell. From this:
\begin{equation}
  R_\eta = \frac{E(3\times7)}{E(7\times7)},
  \qquad
  R_\phi = \frac{E(3\times3)}{E(3\times7)}.
  \label{eq:retarphi}
\end{equation}

% ---------- 3.6 Assumptions ------------------------------------------------
\subsection{Assumptions}
\label{sec:method:assumptions}

The closure test operates under the following assumptions:

\begin{enumerate}
  \item \textbf{Known distortion model.} The pseudo-data are generated from
    the same MC sample using a fully known and reproducible transformation.
    In real data, the data--MC differences are unknown and may not follow the
    assumed bias pattern.
  \item \textbf{Same-sample derivation and application.} The corrections are
    derived and applied using the same events. In practice, one would derive
    corrections from a control sample and apply them to a signal sample. The
    closure test does not assess transfer bias between different event
    populations.
  \item \textbf{Energy conservation.} Both the distortion and the correction
    preserve the total cluster energy. This avoids introducing energy scale
    biases but means the method is insensitive to overall normalisation
    differences.
  \item \textbf{$|\eta|$-binned corrections.} The corrections are derived
    independently in each $|\eta|$~bin. This assumes that the correction
    parameters vary primarily with $|\eta|$ (due to detector geometry changes)
    and are approximately constant within each bin.
  \item \textbf{Additive correction in fraction space.} Both methods apply
    additive corrections to the cell energy fractions, which is the natural
    parameterisation for the cell-energy reweighting method and assumes that
    the data--MC differences can be described as shifts in the normalised
    energy distribution.
\end{enumerate}

% ============================================================================
\section{Adaptations from the Reference Method}
\label{sec:adaptations}
% ============================================================================

The cell-energy reweighting method implemented in this closure test follows
the approach described in ATL-COM-PHYS-2021-640~(Section~5), with several
adaptations necessitated by differences in the input data format between the
Run~2 ntuples used in the original method and the MC23e ntuples used in this
analysis. This section documents these differences and justifies the choices
made.

\subsection{Input data format}
\label{sec:adaptations:input}

The original method operates on \emph{matched} data--MC event pairs, where
each MC photon is paired with a corresponding data photon from the same
$Z\to\ell\ell\gamma$ event. These matched ntuples contain branches for both
the MC and data cell energies, per-cell $\eta$ positions, and per-cell $\phi$
positions for the $7\times11$, $3\times7$, and $5\times5$ cluster windows in
Layers~1, 2, and~3. The photons have already passed tight identification and
isolation requirements, and the matching procedure additionally ensures
high-quality clusters.

The MC23e ntuples used in this analysis contain \emph{unmatched} MC photon
candidates from the NTupleMaker framework. Each event stores the $7\times11$
Layer-2 cell energies as a vector of 77 double-precision values and the
cluster barycentre pseudorapidity $\eta_2$, but does \emph{not} store
per-cell $\eta$ or $\phi$ coordinates. The photon candidates have not
undergone the same level of pre-selection as the matched ntuples.

These format differences require the adaptations described below.

\subsection{Cluster quality selection}
\label{sec:adaptations:quality}

Since the MC23e ntuples contain all photon candidates without prior matching
or tight pre-selection, the sample includes pathological clusters that
would not appear in the original method's input. Two quality requirements
are applied:

\paragraph{Central cell dominance.}
Events are required to have the central cell (cell index~38, corresponding to
$(\eta_\text{row},\phi_\text{col})=(3,5)$) as the highest-energy cell in the
cluster. This removes mis-centred clusters where the shower maximum falls
outside the nominal cluster centre, producing unreliable energy fractions. In
the original method, this is handled implicitly by the event matching and
selection procedure.

\paragraph{Negative energy clamping.}
Cell energies arising from electronic noise (pedestal subtraction) can be
negative. In the original method, this has negligible impact because
the matched ntuples contain only high-quality photons. In the MC23e ntuples,
however, approximately 2.7\% of events have sufficiently large negative cell
contributions that the sum of cell energies becomes small or negative,
producing cell energy fractions $f_k > 1$ or $f_k < 0$. These events
generate pathological corrections that dominate the \texttt{TProfile} means
and can reverse the sign of the flat cell reweighting corrections. Negative cell energies
are therefore clamped to zero before computing energy fractions.

Neither of these requirements is present in the original method because its
input data has already been cleaned by the matching and selection procedure.

\subsection{Conversion type binning}
\label{sec:adaptations:conversion}

The original method derives and applies separate corrections for three photon
categories: inclusive, unconverted, and converted. This accounts for the
different lateral shower profiles of converted and unconverted photons.

The current analysis derives inclusive corrections only, without separation by
conversion type. This simplification is appropriate for the closure test,
where the pseudo-data is generated from the same MC sample and therefore has
the same conversion fraction by construction. For application to real data,
conversion-type binning should be restored.

\subsection{Numerical precision}
\label{sec:adaptations:precision}

The original method uses single-precision floating-point arithmetic
(\texttt{float}) for cell energies, consistent with the Run~2 ntuple format.
The MC23e ntuples store cell energies as double-precision (\texttt{double}).
All calculations in this analysis are performed in double precision. This
difference has negligible impact on the results.

\subsection{Summary of adaptations}
\label{sec:adaptations:summary}

\Cref{tab:adaptations} summarises the differences between the original
method and the current analysis.

\begin{table}[htbp]
\centering
\caption{Summary of adaptations from the reference implementation (ATL-COM-PHYS-2021-640).
  Items marked with $^\dagger$ are specific to the closure test and should be
  revisited for application to real data.}
\label{tab:adaptations}
\small
\begin{tabular}{lll}
\toprule
\textbf{Aspect} & \textbf{Reference method} & \textbf{This analysis} \\
\midrule
Input format          & Matched data--MC (\texttt{float}) & Unmatched MC (\texttt{double}) \\
Pre-selection         & Tight ID + isolation + matching     & Central cell + energy clamping \\
Conversion binning    & Inc / NCon / Con                    & Inclusive only$^\dagger$ \\
TProfile fraction bins & 12 bins in $[-1, 1]$              & Same \\
$|\eta|$ bins         & 14~bins, $[0, 2.4]$                & Same \\
Energy sub-windows    & $E(\Delta\eta, \Delta\phi)$ function & Same \\
Correction formula    & $\langle f^\text{data} {-} f^\text{MC} | f^\text{MC}\rangle$ & Same \\
Correction application & \texttt{TProfile::Interpolate}    & Same \\
\bottomrule
\end{tabular}
\end{table}

% ============================================================================
\section{Validation}
\label{sec:validation}
% ============================================================================

\subsection{Closure metric}
\label{sec:validation:metric}

The primary quantitative measure of closure is the $\chi^2/n_\text{df}$
between the normalised shower shape distributions of the corrected MC and the
pseudo-data target:
\begin{equation}
  \chi^2 = \sum_{b=1}^{N_\text{bins}}
    \frac{\bigl(h_\text{corr}^{(b)} - h_\text{PS}^{(b)}\bigr)^2}
         {h_\text{PS}^{(b)}},
  \label{eq:chi2}
\end{equation}
where $h_\text{corr}^{(b)}$ and $h_\text{PS}^{(b)}$ are the normalised bin
contents of the corrected MC and pseudo-data histograms, respectively, and the
sum runs over all bins with non-zero pseudo-data content. The number of
degrees of freedom $n_\text{df}$ equals the number of non-empty bins. Each
shower shape variable is binned in 25~bins
($R_\eta\in[0.90,1.0]$;
$R_\phi\in[0.85,1.0]$).

A smaller $\chi^2/n_\text{df}$ indicates better agreement. The uncorrected MC
value is also computed as a baseline.

\subsection{Multi-level distortion scan}
\label{sec:validation:scan}

The closure test is performed at three distortion levels
($\delta=0.01,\,0.03,\,0.05$) to assess the linearity and robustness of the
methods. A method that achieves good closure across different distortion
magnitudes demonstrates that it can handle a range of data--MC discrepancy
sizes.

\subsection{Full-statistics validation}
\label{sec:validation:stats}

The validation uses the full MC23e sample of 9.6~million events, of which
approximately 7.8~million pass the quality selection (81\%). The high
statistics ensure that the $\chi^2$ metric is not dominated by statistical
fluctuations and provides a sensitive probe of systematic residuals.

% ============================================================================
\section{Results}
\label{sec:results}
% ============================================================================

\subsection{Representative closure plots at 5\% distortion}
\label{sec:results:representative}

\Cref{fig:5pct_eta0} shows the closure plots for both shower shape
variables in the central barrel region ($|\eta|\in[0.00,\,0.20)$) at the
5\% distortion level. Each plot uses a two-panel format: the upper panel
displays the normalised distributions overlaid for the pseudo-data (black
markers), the uncorrected MC (solid red line), Method~1 reweighted MC (dashed
green line), and Method~2 reweighted MC (solid blue line). The lower panel
shows the ratio of each distribution to the pseudo-data target, with a
horizontal dashed grey line at unity marking perfect agreement.

At 5\% distortion, the uncorrected MC is visibly displaced from the
pseudo-data. Method~1 reduces the discrepancy substantially, and Method~2
brings the ratio to within~1\% of unity across the full range. This pattern
is consistent across both variables.

\begin{figure}[htbp]
  \centering
  \begin{subfigure}[t]{0.48\textwidth}
    \includegraphics[page=1,width=\textwidth]{../output/closure_test_5pct/closure_reta.pdf}
    \caption{$R_\eta$}
    \label{fig:5pct_eta0_reta}
  \end{subfigure}\hfill
  \begin{subfigure}[t]{0.48\textwidth}
    \includegraphics[page=1,width=\textwidth]{../output/closure_test_5pct/closure_rphi.pdf}
    \caption{$R_\phi$}
    \label{fig:5pct_eta0_rphi}
  \end{subfigure}
  \caption{%
    Closure plots for $R_\eta$~(a) and $R_\phi$~(b)
    at 5\% distortion in the central barrel ($|\eta|\in[0.00,\,0.20)$).
    Upper panels: normalised distributions of pseudo-data (markers),
    uncorrected MC (red), Method~1 (green dashed), and Method~2 (blue).
    Lower panels: ratio to pseudo-data.
  }
  \label{fig:5pct_eta0}
\end{figure}

\Cref{fig:5pct_eta13} shows the same two variables at 5\% distortion in
the forward endcap ($|\eta|\in[2.20,\,2.40)$), demonstrating that the
correction achieves comparable closure at high pseudorapidity where the
detector geometry differs from the central barrel. The uncorrected MC
$\chi^2/n_\text{df}$ is substantially larger at high $|\eta|$ (reflecting
different shower development), but Method~2 still achieves a reduction by
approximately a factor of~20.

\begin{figure}[htbp]
  \centering
  \begin{subfigure}[t]{0.48\textwidth}
    \includegraphics[page=13,width=\textwidth]{../output/closure_test_5pct/closure_reta.pdf}
    \caption{$R_\eta$}
  \end{subfigure}\hfill
  \begin{subfigure}[t]{0.48\textwidth}
    \includegraphics[page=13,width=\textwidth]{../output/closure_test_5pct/closure_rphi.pdf}
    \caption{$R_\phi$}
  \end{subfigure}
  \caption{%
    Closure plots at 5\% distortion in the forward endcap
    ($|\eta|\in[2.20,\,2.40)$). Same layout as~\cref{fig:5pct_eta0}.
  }
  \label{fig:5pct_eta13}
\end{figure}

\Cref{fig:distortion_scan} compares $R_\eta$ at 1\%, 3\%, and 5\%
distortion in the central barrel, illustrating how the correction performance
scales with the distortion magnitude. At 1\%, the data--MC differences are
very small and all distributions are nearly indistinguishable by eye. At 3\%,
the displacement of the uncorrected MC becomes visible and both methods
provide good correction. At 5\%, the discrepancy is clearly apparent and the
superiority of Method~2 over Method~1 is most evident.

\begin{figure}[htbp]
  \centering
  \begin{subfigure}[t]{0.32\textwidth}
    \includegraphics[page=1,width=\textwidth]{../output/closure_test_1pct/closure_reta.pdf}
    \caption{$\delta = 0.01$ (1\%)}
    \label{fig:scan_1pct}
  \end{subfigure}\hfill
  \begin{subfigure}[t]{0.32\textwidth}
    \includegraphics[page=1,width=\textwidth]{../output/closure_test_3pct/closure_reta.pdf}
    \caption{$\delta = 0.03$ (3\%)}
    \label{fig:scan_3pct}
  \end{subfigure}\hfill
  \begin{subfigure}[t]{0.32\textwidth}
    \includegraphics[page=1,width=\textwidth]{../output/closure_test_5pct/closure_reta.pdf}
    \caption{$\delta = 0.05$ (5\%)}
    \label{fig:scan_5pct}
  \end{subfigure}
  \caption{%
    Distortion level comparison for $R_\eta$ in the central barrel
    ($|\eta|\in[0.00,\,0.20)$) at 1\%~(a), 3\%~(b), and 5\%~(c)
    distortion. Same layout as~\cref{fig:5pct_eta0}.
  }
  \label{fig:distortion_scan}
\end{figure}

\clearpage

% ---------- 5.2 Closure metrics tables -------------------------------------
\subsection{Closure metrics}
\label{sec:results:chi2}

\Cref{tab:chi2_1pct,tab:chi2_3pct,tab:chi2_5pct} present the
$\chi^2/n_\text{df}$ values for both shower shape variables, in all
populated $|\eta|$ bins, for the three distortion levels. The columns show
the uncorrected MC (``MC''), Method~1 (``M1''), and Method~2 (``M2'').

\begin{table}[htbp]
\centering
\caption{%
  $\chi^2/n_\text{df}$ at 1\% distortion ($\delta=0.01$) for all shower shape
  variables across $|\eta|$ bins.
}
\label{tab:chi2_1pct}
\sisetup{
  table-format=1.4,
  table-number-alignment=center
}
\small
\begin{tabular}{
  c
  S S S
  S S S
}
\toprule
& \multicolumn{3}{c}{$R_\eta$}
& \multicolumn{3}{c}{$R_\phi$} \\
\cmidrule(lr){2-4}\cmidrule(lr){5-7}
{$|\eta|$ bin}
& {MC} & {M1} & {M2}
& {MC} & {M1} & {M2} \\
\midrule
0  & 0.0003 & 0.0011 & 0.0001 & 0.0002 & 0.0004 & 0.0000 \\
1  & 0.0003 & 0.0009 & 0.0001 & 0.0002 & 0.0004 & 0.0000 \\
2  & 0.0003 & 0.0010 & 0.0001 & 0.0002 & 0.0004 & 0.0000 \\
3  & 0.0003 & 0.0009 & 0.0001 & 0.0002 & 0.0003 & 0.0000 \\
4  & 0.0003 & 0.0008 & 0.0001 & 0.0002 & 0.0004 & 0.0000 \\
5  & 0.0003 & 0.0008 & 0.0001 & 0.0002 & 0.0003 & 0.0000 \\
6  & 0.0003 & 0.0007 & 0.0001 & 0.0002 & 0.0003 & 0.0000 \\
7  & 0.0005 & 0.0008 & 0.0001 & 0.0003 & 0.0004 & 0.0001 \\
9  & 0.0004 & 0.0013 & 0.0002 & 0.0001 & 0.0003 & 0.0001 \\
10 & 0.0004 & 0.0009 & 0.0001 & 0.0002 & 0.0003 & 0.0000 \\
11 & 0.0006 & 0.0007 & 0.0001 & 0.0003 & 0.0005 & 0.0000 \\
12 & 0.0007 & 0.0006 & 0.0000 & 0.0004 & 0.0004 & 0.0000 \\
13 & 0.0008 & 0.0006 & 0.0001 & 0.0005 & 0.0004 & 0.0000 \\
\bottomrule
\end{tabular}
\end{table}

\begin{table}[htbp]
\centering
\caption{%
  $\chi^2/n_\text{df}$ at 3\% distortion ($\delta=0.03$) for all shower shape
  variables across $|\eta|$ bins.
}
\label{tab:chi2_3pct}
\sisetup{
  table-format=1.4,
  table-number-alignment=center
}
\small
\begin{tabular}{
  c
  S S S
  S S S
}
\toprule
& \multicolumn{3}{c}{$R_\eta$}
& \multicolumn{3}{c}{$R_\phi$} \\
\cmidrule(lr){2-4}\cmidrule(lr){5-7}
{$|\eta|$ bin}
& {MC} & {M1} & {M2}
& {MC} & {M1} & {M2} \\
\midrule
0  & 0.0023 & 0.0082 & 0.0011 & 0.0017 & 0.0031 & 0.0004 \\
1  & 0.0024 & 0.0073 & 0.0009 & 0.0017 & 0.0032 & 0.0004 \\
2  & 0.0025 & 0.0073 & 0.0009 & 0.0019 & 0.0032 & 0.0004 \\
3  & 0.0024 & 0.0070 & 0.0009 & 0.0017 & 0.0029 & 0.0004 \\
4  & 0.0029 & 0.0062 & 0.0007 & 0.0016 & 0.0033 & 0.0004 \\
5  & 0.0029 & 0.0061 & 0.0006 & 0.0017 & 0.0027 & 0.0003 \\
6  & 0.0028 & 0.0062 & 0.0007 & 0.0016 & 0.0023 & 0.0003 \\
7  & 0.0031 & 0.0054 & 0.0007 & 0.0020 & 0.0025 & 0.0004 \\
9  & 0.0032 & 0.0098 & 0.0014 & 0.0007 & 0.0023 & 0.0003 \\
10 & 0.0035 & 0.0068 & 0.0007 & 0.0017 & 0.0026 & 0.0003 \\
11 & 0.0054 & 0.0060 & 0.0004 & 0.0029 & 0.0040 & 0.0004 \\
12 & 0.0063 & 0.0051 & 0.0003 & 0.0038 & 0.0032 & 0.0003 \\
13 & 0.0074 & 0.0049 & 0.0004 & 0.0042 & 0.0027 & 0.0003 \\
\bottomrule
\end{tabular}
\end{table}

\begin{table}[htbp]
\centering
\caption{%
  $\chi^2/n_\text{df}$ at 5\% distortion ($\delta=0.05$) for all shower shape
  variables across $|\eta|$ bins.
}
\label{tab:chi2_5pct}
\sisetup{
  table-format=1.4,
  table-number-alignment=center
}
\small
\begin{tabular}{
  c
  S S S
  S S S
}
\toprule
& \multicolumn{3}{c}{$R_\eta$}
& \multicolumn{3}{c}{$R_\phi$} \\
\cmidrule(lr){2-4}\cmidrule(lr){5-7}
{$|\eta|$ bin}
& {MC} & {M1} & {M2}
& {MC} & {M1} & {M2} \\
\midrule
0  & 0.0066 & 0.0204 & 0.0028 & 0.0047 & 0.0083 & 0.0012 \\
1  & 0.0068 & 0.0184 & 0.0024 & 0.0049 & 0.0082 & 0.0011 \\
2  & 0.0070 & 0.0185 & 0.0024 & 0.0052 & 0.0083 & 0.0011 \\
3  & 0.0067 & 0.0175 & 0.0023 & 0.0048 & 0.0075 & 0.0010 \\
4  & 0.0085 & 0.0158 & 0.0017 & 0.0047 & 0.0087 & 0.0011 \\
5  & 0.0082 & 0.0157 & 0.0018 & 0.0046 & 0.0070 & 0.0009 \\
6  & 0.0081 & 0.0152 & 0.0018 & 0.0046 & 0.0059 & 0.0009 \\
7  & 0.0084 & 0.0139 & 0.0016 & 0.0049 & 0.0067 & 0.0010 \\
9  & 0.0087 & 0.0249 & 0.0032 & 0.0019 & 0.0059 & 0.0007 \\
10 & 0.0101 & 0.0174 & 0.0017 & 0.0047 & 0.0068 & 0.0008 \\
11 & 0.0153 & 0.0156 & 0.0011 & 0.0085 & 0.0101 & 0.0009 \\
12 & 0.0183 & 0.0131 & 0.0009 & 0.0109 & 0.0080 & 0.0009 \\
13 & 0.0213 & 0.0128 & 0.0009 & 0.0117 & 0.0072 & 0.0008 \\
\bottomrule
\end{tabular}
\end{table}

\clearpage

% ---------- 5.3 Interpretation ---------------------------------------------
\subsection{Interpretation}
\label{sec:results:interpretation}

Several systematic features are observed in the closure metrics:

\paragraph{Method~2 achieves substantial closure; Method~1 does not.}
In all $|\eta|$ bins, at all distortion levels, Method~2
reduces the $\chi^2/n_\text{df}$ by a large factor relative to the
uncorrected MC. Method~1 (flat cell reweighting, §5.1), however, generally increases the
$\chi^2$ for $R_\eta$ and $R_\phi$, particularly in the barrel
($|\eta|<1.37$), demonstrating that a single additive constant cannot
capture the structure of the broadening bias.

\paragraph{Method~2 dramatically outperforms Method~1.}
The energy-dependent cell reweighting achieves $\chi^2/n_\text{df}$ values
typically one to two orders of magnitude smaller than the flat cell
reweighting. For example, at $\delta=0.05$ in $|\eta|$ bin~0, the $R_\eta$
values are: MC~=~0.0066, M1~=~0.0204, M2~=~0.0028. The flat cell
reweighting actually increases the $\chi^2$ well above the uncorrected MC in
most barrel bins, while Method~2 consistently reduces it. The contrast is
most dramatic for $R_\eta$, which is the variable most sensitive to the
lateral energy distribution.

\paragraph{The residual $\chi^2$ scales sub-linearly with distortion level.}
The uncorrected MC $\chi^2$ scales approximately quadratically with $\delta$
(as expected for a squared difference), but the corrected $\chi^2$ for
Method~2 scales more slowly, indicating that the method captures the dominant
structure of the distortion even at larger amplitudes.

\paragraph{The method is uniform across $|\eta|$ bins.}
While the barrel bins ($|\eta|<1.37$) and endcap bins ($|\eta|>1.52$) have
different detector geometries and cell sizes, the correction achieves
comparable $\chi^2$ reduction factors across all bins. The forward endcap bins
($|\eta|>1.8$) show somewhat larger uncorrected $\chi^2$ values but also
correspondingly larger absolute corrections.

\paragraph{Transition-region bins show slightly degraded performance.}
The $|\eta|$ bins adjacent to the barrel--endcap gap (bin~7 at
$[1.30,1.37)$ and bin~9 at $[1.52,1.60)$) show slightly degraded
performance. These bins have unusual detector geometry and lower statistics.
Nevertheless, Method~2 still achieves significant $\chi^2$ reduction in these
bins, while Method~1 increases the $\chi^2$ here as well.

\paragraph{Consistency across shower shape variables.}
Both variables ($R_\eta$, $R_\phi$) show similar patterns
of improvement, confirming that the cell-level correction simultaneously and
consistently corrects all derived shower shape quantities.

\subsection{Summary of performance}
\label{sec:results:summary}

Averaging across $|\eta|$ bins at $\delta=0.05$:
\begin{itemize}
  \item \textbf{Method~1} generally increases the $R_\eta$ and $R_\phi$
    $\chi^2/n_\text{df}$ relative to uncorrected MC in barrel bins; only in
    the most forward endcap bins ($|\eta|>2.0$) does it provide modest
    improvement.
  \item \textbf{Method~2} reduces the $R_\eta$ $\chi^2/n_\text{df}$ by
    approximately 60--95\% relative to uncorrected MC.
  \item Similar patterns are observed for $R_\phi$.
\end{itemize}

% ============================================================================
\section{Conclusions}
\label{sec:conclusions}
% ============================================================================

A closure test of the cell-energy reweighting method has been performed using
MC23e $Z\to ee\gamma$ simulated events, validating the procedure for
correcting photon shower shape variables in the ATLAS electromagnetic
calorimeter.

The test demonstrates that the energy-dependent cell reweighting (Method~2, §5.2)
successfully recovers the target pseudo-data distributions across all
13~populated $|\eta|$ bins and three distortion levels. The flat cell reweighting
(Method~1, §5.1) generally increases the $\chi^2$ in barrel bins,
demonstrating that a single additive constant cannot capture the structure of
the broadening bias. Method~2 achieves consistently superior performance,
reducing the $\chi^2/n_\text{df}$ by approximately 60--95\% relative to
uncorrected MC.

The method is robust across the full pseudorapidity range of the
electromagnetic calorimeter ($|\eta|<2.4$, excluding the transition region
$1.37<|\eta|<1.52$), across three different distortion magnitudes (1\%, 3\%,
5\%), and for both shower shape variables tested ($R_\eta$, $R_\phi$).

These results confirm that the cell-energy reweighting method is a viable
approach for correcting data--MC shower shape discrepancies and that the
energy-dependent cell reweighting (\S5.2) provides a significant improvement over the
flat cell reweighting (\S5.1) baseline. The closure test provides the necessary validation to
proceed with the application of the method to real ATLAS data, where
corrections would be derived by comparing data and MC in a suitable control
region.

% ============================================================================
\clearpage
\appendix
\section{Complete Closure Plots}
\label{app:plots}
% ============================================================================

This appendix presents the full set of closure plots for all shower shape
variables, all $|\eta|$~bins, and all distortion levels. For each $|\eta|$
bin, the two shower shape variables ($R_\eta$, $R_\phi$)
are shown side by side. Each plot uses the same two-panel format described in
\cref{sec:results:representative}: an upper panel overlaying the normalised
distributions (pseudo-data, MC, Method~1, Method~2) and a lower panel showing
the ratio to pseudo-data.

Bin~8 ($|\eta|\in[1.37,1.52)$) corresponds to the barrel--endcap transition
region and is empty.

% --- Helper command for one eta-bin row (2 subfigs) -------------------------
\newcommand{\closurerow}[3]{%
  % #1 = distortion dir (e.g. closure_test_5pct)
  % #2 = page number in the multi-page PDF
  % #3 = eta label (e.g. $|\eta|\in[0.00,\,0.20)$)
  \begin{figure}[htbp]
    \centering
    \begin{subfigure}[t]{0.48\textwidth}
      \includegraphics[page=#2,width=\textwidth]{../output/#1/closure_reta.pdf}
      \caption{$R_\eta$}
    \end{subfigure}\hfill
    \begin{subfigure}[t]{0.48\textwidth}
      \includegraphics[page=#2,width=\textwidth]{../output/#1/closure_rphi.pdf}
      \caption{$R_\phi$}
    \end{subfigure}
    \caption{Closure plots for #3.}
  \end{figure}
}

% ---------- A.1 5% distortion ----------------------------------------------
\subsection{5\% distortion ($\delta = 0.05$)}
\label{app:5pct}

\closurerow{closure_test_5pct}{1}{$|\eta|\in[0.00,\,0.20)$, $\delta=0.05$}
\closurerow{closure_test_5pct}{2}{$|\eta|\in[0.20,\,0.40)$, $\delta=0.05$}
\closurerow{closure_test_5pct}{3}{$|\eta|\in[0.40,\,0.60)$, $\delta=0.05$}
\closurerow{closure_test_5pct}{4}{$|\eta|\in[0.60,\,0.80)$, $\delta=0.05$}
\closurerow{closure_test_5pct}{5}{$|\eta|\in[0.80,\,1.00)$, $\delta=0.05$}
\closurerow{closure_test_5pct}{6}{$|\eta|\in[1.00,\,1.20)$, $\delta=0.05$}
\closurerow{closure_test_5pct}{7}{$|\eta|\in[1.20,\,1.30)$, $\delta=0.05$}
\closurerow{closure_test_5pct}{8}{$|\eta|\in[1.30,\,1.37)$, $\delta=0.05$}
\closurerow{closure_test_5pct}{9}{$|\eta|\in[1.52,\,1.60)$, $\delta=0.05$}
\closurerow{closure_test_5pct}{10}{$|\eta|\in[1.60,\,1.80)$, $\delta=0.05$}
\closurerow{closure_test_5pct}{11}{$|\eta|\in[1.80,\,2.00)$, $\delta=0.05$}
\closurerow{closure_test_5pct}{12}{$|\eta|\in[2.00,\,2.20)$, $\delta=0.05$}
\closurerow{closure_test_5pct}{13}{$|\eta|\in[2.20,\,2.40)$, $\delta=0.05$}

\clearpage
% ---------- A.2 3% distortion ----------------------------------------------
\subsection{3\% distortion ($\delta = 0.03$)}
\label{app:3pct}

\closurerow{closure_test_3pct}{1}{$|\eta|\in[0.00,\,0.20)$, $\delta=0.03$}
\closurerow{closure_test_3pct}{2}{$|\eta|\in[0.20,\,0.40)$, $\delta=0.03$}
\closurerow{closure_test_3pct}{3}{$|\eta|\in[0.40,\,0.60)$, $\delta=0.03$}
\closurerow{closure_test_3pct}{4}{$|\eta|\in[0.60,\,0.80)$, $\delta=0.03$}
\closurerow{closure_test_3pct}{5}{$|\eta|\in[0.80,\,1.00)$, $\delta=0.03$}
\closurerow{closure_test_3pct}{6}{$|\eta|\in[1.00,\,1.20)$, $\delta=0.03$}
\closurerow{closure_test_3pct}{7}{$|\eta|\in[1.20,\,1.30)$, $\delta=0.03$}
\closurerow{closure_test_3pct}{8}{$|\eta|\in[1.30,\,1.37)$, $\delta=0.03$}
\closurerow{closure_test_3pct}{9}{$|\eta|\in[1.52,\,1.60)$, $\delta=0.03$}
\closurerow{closure_test_3pct}{10}{$|\eta|\in[1.60,\,1.80)$, $\delta=0.03$}
\closurerow{closure_test_3pct}{11}{$|\eta|\in[1.80,\,2.00)$, $\delta=0.03$}
\closurerow{closure_test_3pct}{12}{$|\eta|\in[2.00,\,2.20)$, $\delta=0.03$}
\closurerow{closure_test_3pct}{13}{$|\eta|\in[2.20,\,2.40)$, $\delta=0.03$}

\clearpage
% ---------- A.3 1% distortion ----------------------------------------------
\subsection{1\% distortion ($\delta = 0.01$)}
\label{app:1pct}

\closurerow{closure_test_1pct}{1}{$|\eta|\in[0.00,\,0.20)$, $\delta=0.01$}
\closurerow{closure_test_1pct}{2}{$|\eta|\in[0.20,\,0.40)$, $\delta=0.01$}
\closurerow{closure_test_1pct}{3}{$|\eta|\in[0.40,\,0.60)$, $\delta=0.01$}
\closurerow{closure_test_1pct}{4}{$|\eta|\in[0.60,\,0.80)$, $\delta=0.01$}
\closurerow{closure_test_1pct}{5}{$|\eta|\in[0.80,\,1.00)$, $\delta=0.01$}
\closurerow{closure_test_1pct}{6}{$|\eta|\in[1.00,\,1.20)$, $\delta=0.01$}
\closurerow{closure_test_1pct}{7}{$|\eta|\in[1.20,\,1.30)$, $\delta=0.01$}
\closurerow{closure_test_1pct}{8}{$|\eta|\in[1.30,\,1.37)$, $\delta=0.01$}
\closurerow{closure_test_1pct}{9}{$|\eta|\in[1.52,\,1.60)$, $\delta=0.01$}
\closurerow{closure_test_1pct}{10}{$|\eta|\in[1.60,\,1.80)$, $\delta=0.01$}
\closurerow{closure_test_1pct}{11}{$|\eta|\in[1.80,\,2.00)$, $\delta=0.01$}
\closurerow{closure_test_1pct}{12}{$|\eta|\in[2.00,\,2.20)$, $\delta=0.01$}
\closurerow{closure_test_1pct}{13}{$|\eta|\in[2.20,\,2.40)$, $\delta=0.01$}

\end{document}
